% TODO make hw stylesheet
\documentclass[10pt]{scrartcl}
\usepackage[a4paper, total={6in, 8in}]{geometry}
\usepackage[usenames,dvipsnames,svgnames]{xcolor}
\usepackage[shortlabels]{enumitem}
\usepackage[framemethod=TikZ]{mdframed}
\usepackage{amsmath,amssymb,amsthm}
\usepackage[colorlinks]{hyperref}
\usepackage{multicol}
\usepackage{tabularx}

\hypersetup{
	colorlinks=true,
	linkcolor=blue,
	filecolor=magenta,      
	urlcolor=magenta,
	pdftitle={MAT 322 HW}
}

\usepackage{microtype}
\usepackage{mathtools}
\usepackage[headsepline]{scrlayer-scrpage}
\usepackage{thmtools}
\usepackage[textsize=footnotesize]{todonotes}

\mdfdefinestyle{mdbluebox}{roundcorner=10pt,innerbottommargin=9pt,
	linecolor=blue,backgroundcolor=TealBlue!5,}
\declaretheoremstyle[headfont=\sffamily\bfseries\color{MidnightBlue},
mdframed={style=mdbluebox},]{thmbluebox}

\mdfdefinestyle{mdbloobox}{frametitlefont=\bfseries,innerbottommargin=8pt,
	nobreak=true,backgroundcolor=TealBlue!5,linecolor=blue,}
\declaretheoremstyle[headfont=\bfseries\color{MidnightBlue},
mdframed={style=mdbloobox},headpunct={\\[3pt]},postheadspace=0pt,]{thmbloobox}

\mdfdefinestyle{mdredbox}{frametitlefont=\bfseries,innerbottommargin=8pt,
	nobreak=true,backgroundcolor=Salmon!5,linecolor=RawSienna,}
\declaretheoremstyle[headfont=\bfseries\color{RawSienna},
mdframed={style=mdredbox},headpunct={\\[3pt]},postheadspace=0pt,]{thmredbox}

\mdfdefinestyle{mdgreenbox}{linecolor=ForestGreen,backgroundcolor=ForestGreen!5,
	linewidth=2pt,rightline=false,leftline=true,topline=false,bottomline=false,}
\declaretheoremstyle[headfont=\bfseries\sffamily\color{ForestGreen!70!black},
mdframed={style=mdgreenbox},headpunct={ --- },]{thmgreenbox}

\mdfdefinestyle{mdorangebox}{linecolor=RedOrange,backgroundcolor=RedOrange!5,
	linewidth=2pt,rightline=false,leftline=true,topline=false,bottomline=false,}
\declaretheoremstyle[headfont=\bfseries\sffamily\color{RedOrange!70!black},
mdframed={style=mdorangebox},headpunct={ --- },]{thmorangebox}

\mdfdefinestyle{mdblackbox}{linecolor=black,backgroundcolor=RedViolet!5!gray!5,
	linewidth=3pt,rightline=false,leftline=true,topline=false,bottomline=false,}
\declaretheoremstyle[mdframed={style=mdblackbox}]{thmblackbox}

\declaretheoremstyle[spaceabove=6pt,spacebelow=6pt]{basehead}

\declaretheorem[style=basehead,name=Problem]{problem}
\declaretheorem[style=basehead,name=Problem,numbered=no]{problem*}
\declaretheorem[style=thmbloobox,name=Setup]{setup} % for config geo
\declaretheorem[style=thmredbox,name=Example,sibling=problem]{example}
\declaretheorem[style=thmredbox,name=$\bigstar$ Problem,sibling=problem]{niceprob}
\declaretheorem[style=thmbluebox,name=Theorem,sibling=problem]{theorem}
\declaretheorem[style=thmbluebox,name=Corollary,sibling=theorem]{corollary}
\declaretheorem[style=thmbluebox,name=Proposition,sibling=theorem]{prop}
\declaretheorem[style=thmbluebox,name=Lemma,numberwithin=problem]{lemma}
\declaretheorem[style=thmbluebox,name=Theorem,numbered=no]{theorem*}
\declaretheorem[style=thmbluebox,name=Lemma,numbered=no]{lemma*}
\declaretheorem[style=thmgreenbox,name=Claim,numberwithin=problem]{claim}
\declaretheorem[style=thmgreenbox,name=Claim,numbered=no]{claim*}
\declaretheorem[style=thmblackbox,name=Remark,numberwithin=problem]{remark}
\declaretheorem[style=thmblackbox,name=Remark,numbered=no]{remark*}
\declaretheorem[style=thmgreenbox,name=Definition,sibling=theorem]{definition}
\declaretheorem[style=thmgreenbox,name=Definition,numbered=no]{definition*}
\declaretheorem[style=thmorangebox,name=Warning,sibling=theorem]{warning}
\declaretheorem[style=thmorangebox,name=Warning,numbered=no]{warning*}
\declaretheorem[style=thmorangebox,name=Note,numbered=no]{note}
\declaretheorem[style=thmblackbox,name=Exercise,sibling=problem]{exercise}
\declaretheorem[style=thmblackbox,name=Exercise,numbered=no]{exercise*}
\declaretheorem[style=thmblackbox,name=Question,sibling=problem]{question}
\declaretheorem[style=thmblackbox,name=Question,numbered=no]{question*}

\addtolength{\textheight}{3.14cm}
\providecommand{\ol}{\overline}
\providecommand{\ul}{\underline}
\providecommand{\eps}{\varepsilon}
\providecommand{\half}{\frac{1}{2}}
\providecommand{\dang}{\measuredangle} %% Directed angle
\providecommand{\CC}{\mathbb C}
\providecommand{\FF}{\mathbb F}
\providecommand{\NN}{\mathbb N}
\providecommand{\QQ}{\mathbb Q}
\providecommand{\RR}{\mathbb R}
\providecommand{\ZZ}{\mathbb Z}
\providecommand{\dg}{^\circ}
\providecommand{\ii}{\item}
\providecommand{\setdiff}{\smallsetminus}
\providecommand{\alert}[1]{{\sffamily\textbf{\textcolor{blue}{#1}}}}
\providecommand{\inv}{^{-1}}
\providecommand{\mb}[1]{\mathbf{#1}}
\providecommand{\dif}{\;d}

\reversemarginpar

\newenvironment{soln}{\begin{proof}[Solution]}{\end{proof}}
\newenvironment{walkthrough}{\noindent \textbf{Walkthrough.}}{\medskip}
\newlist{walk}{enumerate}{3}
\setlist[walk]{label=\bfseries (\alph*)}

\providecommand{\pow}{\operatorname{Pow}}
\providecommand{\vocab}[1]{{\textbf{\textcolor{ForestGreen}{#1}}}}
\providecommand{\lcm}{\operatorname{lcm}}
\providecommand{\abs}[1]{\left|#1\right|}

\usepackage{asymptote}
\begin{asydef}
	defaultpen(fontsize(10pt));
	unitsize(1cm);
	size(8cm); // set a reasonable default
	usepackage("amsmath");
	usepackage("amssymb");
	settings.tex="pdflatex";
	settings.outformat="pdf";
	// Replacement for olympiad+cse5 which is not standard
	import geometry;
	// recalibrate fill and filldraw for conics
	void filldraw(picture pic = currentpicture, conic g, pen fillpen=defaultpen, pen drawpen=defaultpen)
	{ filldraw(pic, (path) g, fillpen, drawpen); }
	void fill(picture pic = currentpicture, conic g, pen p=defaultpen)
	{ filldraw(pic, (path) g, p); }
	// some geometry
	pair foot(pair P, pair A, pair B) { return foot(triangle(A,B,P).VC); }
	pair orthocenter(pair A, pair B, pair C) { return orthocentercenter(A,B,C); }
	pair centroid(pair A, pair B, pair C) { return (A+B+C)/3; }
	// cse5 abbreviations
	path CP(pair P, pair A) { return circle(P, abs(A-P)); }
	path CR(pair P, real r) { return circle(P, r); }
	pair IP(path p, path q) { return intersectionpoints(p,q)[0]; }
	pair OP(path p, path q) { return intersectionpoints(p,q)[1]; }
	path Line(pair A, pair B, real a=0.6, real b=a) { return (a*(A-B)+A)--(b*(B-A)+B); }
	// cse5 more useful functions
	picture CC() {
		picture p=rotate(0)*currentpicture;
		currentpicture.erase();
		return p;
	}
	pair MP(Label s, pair A, pair B = plain.S, pen p = defaultpen) {
		Label L = s;
		L.s = "$"+s.s+"$";
		label(L, A, B, p);
		return A;
	}
	pair Drawing(Label s = "", pair A, pair B = plain.S, pen p = defaultpen) {
		dot(MP(s, A, B, p), p);
		return A;
	}
	path Drawing(path g, pen p = defaultpen, arrowbar ar = None) {
		draw(g, p, ar);
		return g;
	}
\end{asydef}

\usepackage{mathrsfs}

\ihead{\footnotesize MAT 314 HW08}
\ohead{\footnotesize Last compiled \today}

\title{\vspace{-2em}MAT 314 HW08}
\author{\large Aditya Pahuja}
\date{\large Due 05/08}
\begin{document}
\maketitle

\begin{problem*}1.1.
	Determine the orbit of the polynomial below.
	If the polynomial is symmetric, write it in terms of
	the elementary symmetric functions.
	\begin{enumerate}[(a)]
		\ii $u_1^2u_2 + u_2^2u_3 + u_3^2u_1$, $n=3$
		\ii $(u_1+u_2)(u_2+u_3)(u_3+u_1)$, $n=3$
		\ii $(u_1-u_2)(u_2-u_3)(u_1-u_3)$, $n=3$
		\ii $u_1^3u_2 + u_2^3u_3 + u_3^3u_1 -
		u_1u_2^3 - u_2u_3^3 - u_3u_1^3$, $n=3$
		\ii $u_1^3 + u_2^3 + \cdots + u_n^3$
	\end{enumerate}
\end{problem*}

\begin{soln}
	(a) Since cycling $u_1\to u_2\to u_3\to u_1$
	preserves the polynomial, even permutations (of which there are $3$)
	preserve the polynomial. Thus the size of the orbit is $6/3=2$,
	and the two polynomials in the orbit are
	\[u_1^2u_2 + u_2^2u_3 + u_3^2u_1,\quad u_1u_2^2 + u_2u_3^2 + u_3u_1^2\]
	with the latter obtained by the permutation $(1\;2)$.
	
	(b) This is symmetric: expanding shows that it's equal to
	\begin{align*}
		u_1^2u_2 + u_2^2u_3 + u_3^2u_1 + u_1u_2^2 + u_2u_3^2 + u_3u_1^2
		+ 2u_1u_2u_3 &= (u_1+u_2+u_3)(u_1u_2+u_2u_3+u_3u_1) - u_1u_2u_3\\
		&= \boxed{s_1s_2 - s_3}.
	\end{align*}
	
	(c) Expanding gives
	\[u_1^2u_2 + u_2^2u_3 + u_3^2u_1
	- u_1u_2^2 + u_2u_3^2 + u_3u_1^2.\]
	As with (a) there are two polynomials in the orbit, which are
	\begin{align*}
		u_1^2u_2 + u_2^2u_3 + u_3^2u_1
		- u_1u_2^2 + u_2u_3^2 + u_3u_1^2\\
		u_2^2u_1 + u_3^2u_2 + u_1^2u_3
		- u_2u_1^2 + u_3u_2^2 + u_1u_3^2
	\end{align*}
	
	(d) Same as (a), we can conclude that
	there are only two polynomials in the orbit, and they are
	\begin{align*}
		u_1^3u_2 + u_2^3u_3 + u_3^3u_1 -
		u_1u_2^3 - u_2u_3^3 - u_3u_1^3\\
		u_2^3u_1 + u_3^3u_2 + u_1^3u_3 -
		u_2u_1^3 - u_3u_2^3 - u_1u_3^3
	\end{align*}

	(e) This is symmetric. We first note that
	\[u_1^2 + u_2^2 + \cdots + u_n^2 = s_1^2 - 2s_2\]
	and that
	\[\sum_{sym}u_1^2u_2 = s_1s_2 - ns_3\]
	where the sum is the orbit sum of $u_1^2u_2$. Then,
	\[s_1(u_1^2 +\cdots+ u_n^2) = u_1^3 + \cdots + u_n^3 + \sum_{sym}u_1^2u_2\]
	which simplifies as
	\[u_1^3 + \cdots + u_n^3 = s_1(s_1^2-2s_2) - s_1s_2 + ns_3
	= \boxed{s_1^3 - 3s_1s_2 + ns_3}\]
\end{soln}

\begin{problem*}3.1.
	Let $f$ be a polynomial of degree $n$ with coefficients in $F$
	and let $K$ be a splitting field for $f$ over $F$.
	Prove that $[K:F]$ divides $n!$.
\end{problem*}

\begin{soln}
	First, suppose $f$ is irreducible. Then we prove the claim by induction;
	the result for $n=0$ is trivial ($K=F$).
	In general, if $\alpha$ is a root of $f$, then
	\[[K:F] = [K:K(\alpha)][K(\alpha):F] = n[K:K(\alpha)].\]
	In $K(\alpha)$, $f$ factors as $(x-\alpha)g(x)$; if $f_1$
	is an irreducible factor of $g$ (in $K(\alpha)$),
	then, letting $L_1$ be the splitting field of $f_1$,
	we have by hypothesis that $[L_1:K(\alpha)]$ divides
	$(\deg f_1)!$. Then we can take an irreducible factor $f_2$
	of $g/f_1$ and take the splitting field $L_2/L_1$, and so on,
	giving a chain of extensions
	$K = L_k\supseteq\cdots\supseteq L_1\supseteq K(\alpha)\supseteq K$.
	The degrees of the $f_i$s add up to $n-1$, so that
	\[[K:K(\alpha)] =
	[L_k:L_{k-1}][L_{k-1}:L_{k-2}]\cdots[L_2:L_1][L_1:K(\alpha)]
	= n_k!n_{k-1}!\cdots n_2!n_1!\mid (n-1)!\]
	where the divisibility is a result of the multinomial coefficient
	$\binom{n-1}{n_1,\dots,n_k}$ being an integer; then the factor of $n$
	from $[K:K(\alpha)]$ gives divisibility by $n!$.
	
	For possibly reducible $f$, let $g$ be an irreducible
	factor of $f$ in $F$. Then the splitting field $K_g$ of $g$
	has degree dividing $(\deg g)!$ by the irreducible case
	and $[K:K_g]$ divides $(\deg f/g)! = (\deg f - \deg g)!$ by induction, so
	\[[K:F] \mid (\deg g)!(n - \deg g)!\mid n!\]
	with divisibility coming from $\binom n{\deg g}$ being an integer.
\end{soln}

\begin{problem*}6.1.
	Let $\alpha$ be a complex root of the polynomial $x^3+x+1$
	over $\QQ$, and let $K$ be a splitting field of this polynomial over $\QQ$.
	Is $\sqrt{-31}$ in the field $\QQ(\alpha)$? Is it in $K$?
\end{problem*}

\begin{soln}
	$\sqrt{-31}$ is not in $\QQ(\alpha)$: if it were, then we would have
	\[3 = [\QQ(\alpha):\QQ] = [\QQ(\alpha):\QQ(\sqrt{-31})]
	[\QQ(\sqrt{-31}):\QQ] = 2[\QQ(\alpha):\QQ(\sqrt{-31})],\]
	which is impossible because the degree of the extension is an integer.
	
	However, $\sqrt{-31}$ is in $K$, because $\sqrt{-31}$ is the square root
	of the discriminant $-4(1) - 27(1) = -31$ of $x^3+x+1$.
\end{soln}

\begin{problem*}7.2.
	Let $K/F$ be a Galois extension such that $G(K/F)\cong C_2\times C_{12}$.
	How many intermediate fields $L$ are there with
	(a) $[L:F] = 4$, (b) $[L:F] = 9$, (c) $G(K/L)\cong C_4$?
\end{problem*}

\begin{soln}
	Because $K/F$ is Galois, $24 = |G(K/F)| = [K:F]$.
	
	(a) $[L:F]=4$ if and only if $[K:L] = 6$.
	Such fields $L$ correspond bijectively to subgroups of $C_2\times C_{12}$
	with order $6$. Noting that all the subgroups are abelian
	(because the group itself is abelian), every order $6$ subgroup
	of $G(K/F)$ must be isomorphic to $C_6$. The order $6$ elements
	in $C_2\times C_{12}$ are
	\[(1, 2), (1, 4), (1, 8), (1, 10), (0, 2), (0, 10),\]
	of which $(1, 2)$ and $(1, 10)$ generate the same subgroup,
	$(1, 4)$ and $(1, 8)$ generate the same subgroup,
	and $(0, 2)$ and $(0,10)$ generate the same subgroup.
	Therefore $C_2\times C_{12}$ has three subgroups of order $6$,
	so there are $\boxed 3$ fields $L$ with the desired property.
	
	(b) $[L:F]=9$ is impossible, since this would imply $[K:L] = 24/9\notin\ZZ$.
	
	(c) Each subgroup isomorphic to $C_4$ corresponds to
	exactly one intermediate field with Galois group isomorphic to $C_4$
	(with the field being the fixed field of the subgroup,
	as per the Galois correspondence).
	We therefore need to count the number of subgroups isomorphic to $C_4$.
	These are generated by elements of order $4$,
	of which there are four:
	\[(0,3), (0, 9), (1, 3), (1, 9).\]
	The first two generate the same subgroup and the latter two
	generate the same subgroup, giving $\boxed{2}$ subgroups of $G(K/F)$
	isomorphic to $C_4$.
\end{soln}

\pagebreak
\begin{problem*}10.3.
	Let $\zeta = \zeta_7$. Determine the degree of
	the following elements over $\QQ$:
	\begin{enumerate}[(a)]
		\ii $\zeta + \zeta^5$,
		\ii $\zeta^3 + \zeta^4$,
		\ii $\zeta^3+\zeta^5+\zeta^6$
	\end{enumerate}
\end{problem*}

\begin{soln}
	(a) The degree is $\boxed 6$.
	
	We know that the subfield generated by $\zeta+\zeta^5$
	is a subfield of $\QQ(\zeta)$, and that the Galois group of
	$\QQ(\zeta)$ is the set of automorphisms $\zeta\mapsto\zeta^k$
	for $k=1,\dots,6$, which we will denote with $\phi_k$.
	This means that $\zeta+\zeta^5$ shares its minimal polynomial with
	\begin{align*}
		\phi_1(\zeta+\zeta^5) &= \zeta+\zeta^5\\
		\phi_2(\zeta+\zeta^5) &= \zeta^2+\zeta^3\\
		\phi_3(\zeta+\zeta^5) &= \zeta^3+\zeta^1\\
		\phi_4(\zeta+\zeta^5) &= \zeta^4+\zeta^6\\
		\phi_5(\zeta+\zeta^5) &= \zeta^5+\zeta^4\\
		\phi_6(\zeta+\zeta^5) &= \zeta^6+\zeta^2,
	\end{align*}
	and no other complex numbers
	(because $\QQ(\zeta)$ is a Galois extension, hence a splitting field).
	Recalling that $\zeta$, \dots, $\zeta^6$ are a basis for $\QQ(\zeta)$,
	we can conclude that $\phi_i(\zeta+\zeta^5)=\phi_j(\zeta+\zeta^5)$
	cannot hold for $i\ne j$, since this would imply linear dependence
	between at most four powers of $\zeta$. This means that
	the minimal polynomial of $\zeta+\zeta^5$ has six distinct roots,
	i.e. the degree of $\zeta+\zeta^5$ is $6$.
	
	(b) The degree is $\boxed{3}$. Letting $x = \zeta^3+\zeta^4$,
	we can see that (noting $\zeta^3\cdot\zeta^4 = 1$)
	\[x^2-2 = \zeta+\zeta^6,\quad x^3-3x = \zeta^2+\zeta^5.\]
	Therefore
	\[x^3+x^2-2x-1 = 1+x+x^2-2+x^3-3x=1+\zeta+\zeta^2+\cdots+\zeta^6=0\]
	and the polynomial on the LHS doesn't factor because
	it has no rational roots (as seen by applying the rational root theorem
	and checking the only two candidates $\pm 1$).
	We have shown that $x$ is the root of an irreducible polynomial
	with degree $3$, which is what we wanted.
	
	(c) The degree is $\boxed2$. If $x = \zeta^3+\zeta^5+\zeta^6$, then
	\[x^2 = \zeta^6 + \zeta^3 + \zeta^5 + 2(\zeta + \zeta^2 + \zeta^4)
	= -2 - x.\]
	Therefore the minimal polynomial of $x$ is $x^2 + x + 2$,
	so $x$ has degree $2$ over $\QQ$.
\end{soln}






















\end{document}